\section{Preparación y validación de datos}

\subsection{Lectura robusta de archivos oficiales}

El software trabaja con archivos Excel o XLSB descargados por el usuario desde fuentes oficiales del DANE. La lectura robusta es un requisito metodológico porque los anexos pueden actualizarse con nuevos periodos, conservar tablas con encabezados complejos o variar ligeramente su estructura visual. El sistema no debe depender del nombre del archivo ni generar meses de forma artificial. Debe identificar los periodos reales presentes en el documento.

La decisión de leer periodos desde encabezados reales tiene una justificación estadística directa: el eje temporal es la base del modelo. Si se generan columnas por conteo, existe riesgo de asignar un mes incorrecto a un dato, desplazar la serie o proyectar desde un último periodo falso. Estos errores no siempre producen fallos visibles en el programa, pero sí contaminan el análisis predictivo.

\subsection{Normalización de periodos}

Los anexos pueden representar periodos con ligeras diferencias textuales, por ejemplo \(2025\_1\), \(2025\_01\) o \(2025\_11\) con espacios al final. El software normaliza estas etiquetas a una forma canónica:

\begin{equation}
    \text{periodo} = \text{año}\_\text{mes},
    \label{eq:periodo_canonico}
\end{equation}

donde el mes se expresa como entero entre 1 y 12. Esta normalización permite ordenar cronológicamente la serie mediante la clave:

\begin{equation}
    k(p) = \left(\text{año}(p), \text{mes}(p)\right).
    \label{eq:period_key}
\end{equation}

Ordenar por texto puede producir errores. Por ejemplo, \(2025\_10\) podría ubicarse antes de \(2025\_2\) si se usa orden lexicográfico. Al convertir año y mes a enteros, el software preserva la secuencia temporal correcta.

\subsection{Validaciones previas a la modelación}

Antes de ajustar modelos, el sistema debe validar la calidad de la serie. Una serie con errores temporales puede generar métricas aparentemente aceptables y, aun así, producir conclusiones inválidas. Las validaciones principales son:

\begin{longtable}{p{0.28\textwidth}p{0.62\textwidth}}
\toprule
\textbf{Validación} & \textbf{Justificación metodológica} \\
\midrule
\endhead
Continuidad temporal & Verifica que no existan meses ausentes entre el primer y último periodo. Un salto temporal altera la distancia entre observaciones. \\
Orden cronológico & Asegura que la serie respete la secuencia pasado-futuro. Sin orden, el backtesting pierde validez. \\
Valores faltantes & Identifica periodos sin índice. Los faltantes pueden sesgar ajustes o reducir el tamaño efectivo de muestra. \\
Duplicados & Evita que un mismo periodo sea contado dos veces. Los duplicados distorsionan pesos implícitos del ajuste. \\
Longitud mínima & Evalúa si existe información suficiente para modelar tendencia y validar desempeño. \\
Outliers extremos & Detecta valores inusuales que podrían corresponder a choques reales o errores de lectura. \\
\bottomrule
\end{longtable}

\subsection{Longitud mínima}

El módulo estadístico recomienda una longitud mínima de 24 observaciones mensuales para modelación. Este umbral no es una ley universal, pero responde a la necesidad de contar con al menos dos años de información para observar tendencia, variabilidad y posibles patrones anuales. El sistema puede emitir advertencia si existen menos de 18 meses y recomendar cautela si hay menos de 24.

La longitud mínima es importante porque varios diagnósticos pierden confiabilidad con muestras pequeñas. Por ejemplo, las pruebas de normalidad de residuos y el cálculo de criterios de información pueden ser inestables cuando hay pocos datos. Además, el backtesting temporal requiere reservar periodos futuros dentro de la muestra histórica; si la serie es demasiado corta, quedan pocas iteraciones de validación.

\subsection{Valores faltantes y tratamiento}

Un valor faltante en una serie ICOCIV no debe rellenarse automáticamente sin justificación. La interpolación puede crear datos que no fueron publicados y alterar la estructura real de la serie. El enfoque metodológico recomendado es reportar los faltantes, impedir la modelación si afectan la continuidad crítica o permitir análisis solo cuando el usuario y el informe reconozcan explícitamente la limitación.

En la práctica del software, la serie usada para regresión se construye a partir de valores numéricos válidos. Si se eliminan valores faltantes, debe verificarse que ello no rompa la continuidad temporal de manera silenciosa. Por eso, la validación de periodos se realiza antes de la modelación.

\subsection{Outliers}

Un outlier es un valor extremo respecto al comportamiento usual de la serie. En índices económicos, no todo outlier es un error. Un aumento fuerte en acero o combustibles puede ser económicamente real. Por tanto, la detección de outliers debe interpretarse como alerta, no como orden automática de eliminación.

Una regla simple para detectar valores extremos en variaciones mensuales es usar puntajes robustos basados en la mediana:

\begin{equation}
    z^{\ast}_t = \frac{x_t - \text{mediana}(x)}{\text{MAD}(x)},
    \label{eq:z_robusto}
\end{equation}

donde \(\text{MAD}\) es la desviación absoluta mediana. Esta estrategia es menos sensible a extremos que la desviación estándar tradicional. En el contexto del software, el resultado debe mostrarse como advertencia documental para que el analista revise si el valor responde a un evento económico justificable.

\subsection{Trazabilidad de la preparación}

La validación de datos no es solo una etapa técnica interna. Debe quedar documentada en el informe porque afecta la confiabilidad del resultado. Un reporte metodológicamente sólido debe indicar número de observaciones, último periodo usado, faltantes, duplicados, continuidad temporal y advertencias relevantes. Así se fortalece la reproducibilidad del análisis y se evita que la proyección aparezca como un resultado aislado sin control de calidad.
